\section{Consequences and further instantiation}
\label{sec:discussion}

\paragraph{A recovery condition attached to original operations.}
The model distinguishes creation of an operation's random input, its first
durable reservation, the acquisition of a retained representation, and the
later availability of the matching decryption material. This distinction
makes exact replay compatible with precise recovery statements: retries and
aggregate writes preserve the original source, while fresh operations obtain
new origins. The guarded characterization explains how records and keys
combine to expose that source. The native executions demonstrate why both
this conjunction and the original event history matter at the persistence
and asynchronous-confirmation boundaries.

\paragraph{Joint root and message instantiation.}
Theorem~\ref{thm:joint-message-leaf} follows the stated finite interface through
confirmation, traffic derivation and the selected first message. Private
Expand-table normalization allows traffic information to depend on
root-derived confirmation tags and supplies the independent hidden input for
the next reduction. Real confirmation and ACK verification, successor-root
disclosure, and disclosure of the next chain seed remain in that continuation.
Exact retries reuse the encryption result from its first private computation.

The connection begins with the root derivation of Equation~\eqref{eq:root},
\[
 p_e=\operatorname{HMAC}_{R_{e-1}}(S_e).
\]
When $R_{e-1}$ is disclosed, the secret contribution occupies HMAC's message
input, the position studied by extractor and swap-PRF
formulations~\cite{hmacswap2023,rfc5869}.
Appendix~\ref{app:joint-root-entry} gives the shared-oracle Extract/Expand
reduction for a uniform combined contribution. Together with the joint
schedule and the specified one-use AEAD game, it bounds the profile's
ideal-contribution score. Section~\ref{sec:contextbound-connection} supplies
the preceding complete-tuple comparison and original-operation adapter.

The executable comparison tracks four identities across failed writes and
exact resumption: the original reservation, the component's challenge-ciphertext
class, the effective hash tuple, and the root-use descriptor. The next-rekey
source profile fixes that descriptor, and the message continuation selects a
fresh allocation. Corollary~\ref{cor:live-core-next-rekey} substitutes the
explicit image/intent comparison cost for the authoritative live core.
The executed-slot selector then associates accepted allocations with the
original reservation pair, up to the signature and framed-hash terms in
Proposition~\ref{prop:executed-source-attribution}. Its signing adapter
preserves original randomizers, exact retries and the declared owner
observations. The native ML-DSA bridge then expresses that signature term through
ordinary hedged unforgeability, an exact-replay coupling and a shared-XOF
query loss (Theorem~\ref{thm:native-mldsa-reservation-bridge}).
The native deterministic P-256 component supplies the corresponding
classical ordinary-signature instantiation.

Proposition~\ref{prop:confirmation-message-auth} follows the same joint
key table through the native confirmation and first-message interfaces.
Its origin ledgers distinguish a fresh authenticated acceptance from an
exact replay or historical consumption status. The joint schedule is charged
once for the union of the confirmation and message events; the remaining
terms account for the MAC rows, complete-core hash collisions and one-key
ciphertext integrity.

Section~\ref{sec:multiple-root-entry} generalizes the entry step to a finite
family of available salts and complete root-information strings. Lifetime
pad-state classes preserve exact HMAC aliases; successive row comparisons
produce an independent table of root entries. A fresh entry in a previously
used class can initialize the marked schedule, while other roots remain
available for later computation. The target's aliases share one schedule
handle, and the complete budgets include public class scheduling and every
other real root computation. A concrete source workload can use this result
by realizing the declared available-data controller.

Section~\ref{sec:sdk-purpose-connection} supplies such a workload in the
SDK's purpose-key API. Its exact information encoding separates purposes,
labels, contexts and policy bindings. The comparison preserves a shared key
entry for equal complete inputs and an independent lifecycle for each
returned owner. It permits all derived-key exports in the joint key-table
experiment. A selected private key additionally supports the stated
application-encryption experiment when all its aliases follow that
experiment's key-use interface.

\paragraph{Matched protocol evaluation.}
The reference corpus fixes reproducible communication opportunities for a
matched Continuity comparison. Such a comparison records the same loss,
reorder and duplication schedule, together with the construction's actual
authentication authority and disclosure objects. For each case, the relevant
observations are the original generation and encapsulation reservations,
acquired record projection, known wrapping keys, public-wire availability,
and reconstructed message set. Communication and persistence costs can then
be measured over declared operation intervals. Holding a reservation fixed
while varying execution time, acknowledgement delivery or aggregate
rewriting tests the provenance distinctions isolated by the present results.

\paragraph{Conclusion.}
Persistent recovery is governed by original secret sources and their
acquisition histories. The paper gives an exact record-exposure semantics,
constructive simulations for private records and dependent immutable reports,
a complete-tuple component connection to the joint first-message continuation,
a finite entry comparison with an owned SDK purpose-key instance, and a mechanized multi-epoch recovery
theorem with concrete separating executions.
These results preserve the meaning of
fresh randomness across durable retries, copies and asynchronous progress.
